\subsection{Properties of modules of fractions}

Throughout this no., A denotes a ring and S a multiplicative subset of A.

Let \((\mathbf{M}_{\alpha}, \phi_{\beta\alpha})\) be a direct system of A-modules; then \((\mathrm{S}^{-1}\mathbf{M}_{\alpha}, \mathrm{S}^{-1}\phi_{\beta\alpha})\) is a direct system of \(S^{-1}A\) -modules and the fact that taking direct limits commutes with tensor products (Algebra, Chapter II, § 6, no. 3, Proposition 7) allows us to define a canonical isomorphism

\[
\lim _ {\longrightarrow} \left(\mathrm{S} ^ {- 1} \mathrm{M} _ {\alpha}\right)\rightarrow \mathrm{S} ^ {- 1} \lim _ {\longrightarrow} \mathrm{M} _ {\alpha}.
\]

Similarly, the fact that taking direct sums commutes with tensor products (Algebra, Chapter II, § 3, no. 7, Proposition 7) allows us to define for every family \((\mathbf{M}_t)_{i\in I}\) of A-modules a canonical isomorphism

\[
\bigoplus_ {i \in I} S ^ {- 1} M _ {i} \rightarrow S ^ {- 1} \bigoplus_ {i \in I} M _ {i}.
\]

Finally we note that, if an A-module M is the sum of a family \((\mathbf{N}_{\iota})_{\iota\in I}\) of submodules, \(S^{-1}M\) is the sum of the family of sub- \(S^{-1}A\) -modules generated by the \(i_{M}^{S}(N_{\iota})\) . Then it follows that, if M is a finitely generated A-module (resp. finitely presented A-module), \(S^{-1}M = S^{-1}A \otimes_{A} M\) is a finitely generated \(S^{-1}A\) -module (resp. finitely presented \(S^{-1}A\) -module).

THEOREM 1. The ring \(\mathbf{S}^{-1}\mathbf{A}\) is a flat A-module (Chapter I, § 2, no. 1, Definition 2).

If \(u: \mathbf{M}' \to \mathbf{M}\) is an injective homomorphism of A-modules, it is necessary to establish that \(\mathbf{S}^{-1}u: \mathbf{S}^{-1}\mathbf{M}' \to \mathbf{S}^{-1}\mathbf{M}\) is injective. Now, if \(m'/s\) ( \(m' \in \mathbf{M}'\) , \(s \in \mathbf{S}\) ) is such that \(u(m')/s = 0\) , this implies the existence of an \(s' \in \mathbf{S}\) such that \(s'u(m') = 0\) (no. 2, Proposition 4) or \(u(s'm') = 0\) ; as \(u\) is injective, it follows that \(s'm' = 0\) , whence \(m'/s = 0\) .

The fact that \(\mathbf{S}^{-1}\mathbf{A}\) is a flat A-module allows us to apply to it the results of Chapter I, § 2. In particular:

\begin{enumerate}
\def\labelenumi{(\arabic{enumi})}
\tightlist
\item
  If M is an A-module and N a submodule of M, \(S^{-1}N\) is canonically identified with a submodule of \(S^{-1}M\) generated by \(i_{\mathrm{M}}^{\mathrm{S}}(\mathrm{N})\) (Chapter I, § 2, no. 3, Remark 2); with this identification, \(S^{-1}(\mathrm{M}/\mathrm{N})\) is identified with \((\mathrm{S}^{-1}\mathrm{M})/(\mathrm{S}^{-1}\mathrm{N})\) and, if P is a second submodule of M, then
\end{enumerate}

\[
\mathrm{S} ^ {- 1} (\mathrm{N} + \mathrm{P}) = \mathrm{S} ^ {- 1} \mathrm{N} + \mathrm{S} ^ {- 1} \mathrm{P}, \quad \mathrm{S} ^ {- 1} (\mathrm{N} \cap \mathrm{P}) = \mathrm{S} ^ {- 1} \mathrm{N} \cap \mathrm{S} ^ {- 1} \mathrm{P}
\]

(Chapter I, § 2, no. 6, Proposition 2).

\begin{enumerate}
\def\labelenumi{(\arabic{enumi})}
\setcounter{enumi}{1}
\tightlist
\item
  If M is a finitely generated A-module, then
\end{enumerate}

\[
\mathbf {S} ^ {- 1} \operatorname{Ann} (\mathbf {M}) = \operatorname{Ann} (\mathbf {S} ^ {- 1} \mathbf {M})\tag{9}
\]

(Chapter I, § 2, no. 10, Corollary 2 to Proposition 12).

PROPOSITION 10. Let M be an A-module. For every submodule N' of the S \(^{-1}\) A-module S \(^{-1}\) M, let \(\phi(N')\) be the inverse image of N' under \(i_{M}^{8}\) . Then:

\begin{enumerate}
\def\labelenumi{(\roman{enumi})}
\item
  \(\mathbf{S}\phi(\mathbf{N}^{\prime}) = \mathbf{N}^{\prime}.\)
\item
  For every submodule \(\mathbf{N}\) of \(\mathbf{M}\) , the submodule \(\phi (\mathrm{S}^{-1}\mathrm{N})\) of \(\mathbf{M}\) consists of those \(m\in \mathbf{M}\) for which there exists \(s\in \mathbf{S}\) such that \(sm\in \mathbf{N}\) .
\item
  \(\phi\) is an isomorphism (for the orderings defined by inclusion) of the set of sub- \(\mathbf{S}^{-1}\mathbf{A}\) -modules of \(\mathbf{S}^{-1}\mathbf{M}\) onto the set of submodules \(\mathbf{Q}\) of \(\mathbf{M}\) which satisfy the following condition:

(MS) If \(sm \in \mathbb{Q}\) , where \(s \in \mathbb{S}\) , \(m \in \mathbb{M}\) , then \(m \in \mathbb{Q}\) .

\end{enumerate}
Obviously \(\mathrm{S}^{-1}\phi(\mathrm{N}')\subset\mathrm{N}'\) ; conversely, if \(n' = m/s \in N'\) , then \(m/1 \in N'\) , hence \(m \in \phi(\mathrm{N}')\) and consequently \(n' \in \mathrm{S}^{-1}(\phi(\mathrm{N'}))\) ; whence (i). For an element \(m \in M\) to be such that \(m \in \phi(\mathrm{S}^{-1}\mathrm{N})\) , it is necessary and sufficient that \(m/1 \in S^{-1}N\) , that is that there exist \(s \in S\) and \(n \in N\) such that \(m/1 = n/s\) ; this means that there exists \(s' \in S\) such that \(s'sm = s'n \in N\) , whence (ii). Finally, the relation \(sm \in \phi(\mathrm{N}')\) is equivalent by definition to \(sm/1 \in N'\) and as s/1 is invertible in \(S^{-1}A\) , this implies \(m/1 \in N'\) , or \(m \in \phi(\mathrm{N}')\) , hence \(\phi(\mathrm{N}')\) satisfies condition (MS); on the other hand, it follows from (ii) that, if N satisfies (MS), then \(\phi(\mathrm{S}^{-1}\mathrm{N}) = N\) , which completes the proof of (iii).

The submodule \(\phi(\mathrm{S}^{-1}\mathrm{N})\) is called the saturation of N in M with respect to S, and the submodules satisfying condition (MS) (and hence equal to their saturations) are said to be saturated with respect to S. The submodule \(\phi(\mathrm{S}^{-1}\mathrm{N})\) is the kernel of the composite homomorphism

\[
\mathrm{M} \xrightarrow {h} \mathrm{M} / \mathrm{N} \xrightarrow {i _ {\mathrm{M} / \mathrm{N}} ^ {\mathrm{S}}} \mathrm{S} ^ {- 1} \mathrm{M} / \mathrm{S} ^ {- 1} \mathrm{N}
\]

where h is the canonical homomorphism, as follows from the commutativity of the diagram

\[
\begin{array}{c c c} \mathrm{M} & \stackrel {{h}} {{\longrightarrow}} & \mathrm{M/N} \\ i _ {\mathrm{M}} ^ {\mathrm{S}} \Big \downarrow & & \Big \downarrow i _ {\mathrm{M/N}} ^ {\mathrm{S}} \\ \mathrm{S} ^ {- 1} \mathrm{M} & \xrightarrow [ \mathrm{S} ^ {- 1 h} ]{} \mathrm{S} ^ {- 1} \mathrm{M/S} ^ {- 1} \mathrm{N} \end{array}
\]

If S is the complement in A of a prime ideal p, \(\phi(S^{-1}N)\) is also called the saturation of N in M with respect to p.

COROLLARY 1. Let \(N_{1}\) , \(N_{2}\) be two submodules of an A-module M. For \(S^{-1}N_{1} \subset S^{-1}N_{2}\) , it is necessary and sufficient that the saturation of \(N_{1}\) with respect to S be contained in that of \(N_{2}\) .

COROLLARY 2. If M is a Noetherian (resp. Artinian) A-module, S \(^{-1}\) M is a Noetherian (resp. Artinian) S \(^{-1}\) A-module. In particular, if the ring A is Noetherian (resp. Artinian), so is the ring S \(^{-1}\) A.

